\bookchapter{Capítulo 5: Deixe o Smartphone Mais Inteligente — Configurações Básicas}

\chapterintro{Dedique alguns minutos às configurações e a experiência melhora enormemente. Essas configurações são como “afinar” o telefone para que ele se adapte melhor aos seus hábitos de uso.}

\booksection{I. Som e Vibração}

Escolha um toque e som de notificação que você goste. Use com flexibilidade o modo silencioso (reuniões) e o modo Não Perturbe (dormir). O mundo fica instantaneamente em paz.

\begin{itemize}
  \item \textbf{Toque}: Escolha um que você consiga ouvir e que não seja constrangedor.
  \item \textbf{Modo de vibração}: No silencioso, a vibração avisa.
  \item \textbf{Modo Não Perturbe}: Pode configurar exceções, como permitir chamadas dos contatos.
  \item \textbf{Separação de volumes}: Volume de mídia, volume de toque e volume do alarme podem ser ajustados separadamente.
\end{itemize}

\booksection{II. Tela e Brilho}

\begin{itemize}
  \item \textbf{Brilho automático}: Deixe o telefone ajustar automaticamente conforme a luz ambiente, sem preocupações.
  \item \textbf{Modo de proteção ocular / Night Shift}: Filtra a luz azul, tornando mais confortável usar o telefone à noite.
  \item \textbf{Modo escuro}: Economiza bateria (especialmente eficaz em telas OLED) e não ofusca à noite. Muitos apps também suportam o modo escuro.
  \item \textbf{Taxa de atualização da tela}: Alguns Androids permitem ajuste. Quanto maior (ex.: 120 Hz), mais fluida é a rolagem.
  \item \textbf{Bloqueio automático}: Configure para 15 ou 30 segundos. Seguro e econômico.
  \item \textbf{Tamanho da fonte}: Acha as letras pequenas? Aumente um pouco.
\end{itemize}

\booksection{III. Papel de Parede, Temas e Área de Trabalho}

Troque por um papel de parede bonito (estático ou dinâmico), o humor muda totalmente. Marcas de smartphones Android geralmente têm lojas de temas que permitem trocar ícones e estilos por completo.

\textbf{Widgets na área de trabalho}: Clima, agenda, controle de música na tela inicial, acesso direto. Tanto iOS quanto Android suportam. Pressione e segure um espaço vazio na área de trabalho para adicionar.

\booksection{IV. Configurações de Produtividade: Transforme o Smartphone em seu “Cérebro Externo”}

Esta seção traz funções das quais “uma vez que você as configura, não consegue mais voltar atrás”. Cada uma delas pode aumentar significativamente sua eficiência diária.

\booksubsection{Modo Foco / Não Perturbe: Recupere a atenção roubada pelas notificações}

\begin{itemize}
  \item \textbf{Ideia central}: Não se trata de "usar menos o telefone", mas de "focar quando é preciso focar, e se conectar quando é preciso se conectar".
  \item \textbf{Sugestões de uso}:\newline · Modo trabalho: permitir apenas contatos "favoritos" e apps essenciais;\newline · Modo sono: ativar automaticamente, permitir apenas alarme;\newline · Exercício / Leitura / Reunião: personalizar conforme a atividade atual.
  \item \textbf{Uso avançado}: Use automação + gatilho de localização/horário para alternar modos automaticamente. Exemplo: o Atalhos do iOS pode ser configurado para "ao chegar na empresa, entrar automaticamente no modo trabalho".
\end{itemize}

\booksubsection{Gerenciamento de Notificações: Faça os pontos vermelhos pararem de “perseguir você”}

\begin{itemize}
  \item \textbf{Desative notificações de marketing}: 90\% das notificações push você realmente não precisa.
  \item \textbf{Ordene por importância}: Priorize comunicações interpessoais (chamadas, SMS, mensagens instantâneas); notícias, compras, redes sociais podem ser definidas como "push silencioso".
  \item \textbf{Notificações agrupadas}: Agrupe notificações semelhantes para evitar que a tela de bloqueio seja inundada por dezenas de mensagens.
\end{itemize}

\booksubsection{Gestos Rápidos: Operações mais ágeis}

\begin{itemize}
  \item \textbf{Captura de tela}: Power + Volume para baixo (iOS tela cheia) / Power + Volume para baixo (Android universal).
  \item \textbf{Captura longa / Captura com rolagem}: Após a captura normal, toque na opção "Captura longa".
  \item \textbf{Escanear com um toque}: Adicione o ícone "Escanear" no Centro de Controle / Configurações rápidas.
  \item \textbf{Toque na parte traseira (iPhone)}: Em [Ajustes] > [Acessibilidade] > [Toque] > [Toque na Parte Traseira], vincule ações como "Captura de tela" ou "Bloquear tela".
  \item \textbf{Despertar por voz}: Siri / Xiao Ai / Xiao Yi / Xiao Bu. Basta chamar para consultar o clima ou definir um alarme.
\end{itemize}

\booksubsection{Teclado: A ferramenta de produtividade mais subestimada}

\begin{itemize}
  \item \textbf{Escolha o teclado certo}:\newline · Sogou / Baidu: vocabulário amplo, atualização rápida de termos populares;\newline · Xunfei: maior taxa de reconhecimento de entrada por voz, ferramenta ideal para anotações em reuniões;\newline · Gboard (Google): simples, sincronização forte entre plataformas, suporta digitação por deslize;\newline · iOS nativo: melhor privacidade, profundamente integrado ao sistema.
  \item \textbf{Opções essenciais}: Correção inteligente, palavras sugeridas, sincronização de dicionário na nuvem, deslizar para cima para inserir números e símbolos.
  \item \textbf{Dica de privacidade}: Cuidado com funções como "sugestão de frases completas" ou "recomendações personalizadas", que podem enviar o conteúdo digitado.
\end{itemize}

\booksubsection{Área de Trabalho e Multitarefa: Fluxo de trabalho mais fluido}

\begin{itemize}
  \item \textbf{Tela dividida / Picture-in-Picture}: Responder mensagens enquanto assiste a vídeos, fazer anotações enquanto navega.
  \item \textbf{Janela flutuante / Mini janela}: WeChat / Douyin suportam "modo mini janela", responda e recolha em seguida.
  \item \textbf{Fixar apps frequentes}: Na Dock (parte inferior da tela), coloque apenas os 4 apps mais frequentes para reduzir o tempo de busca.
\end{itemize}

\begin{tcbinfo}{Checklist do "Primeiro Dia de Trabalho"}
Ao pegar um telefone novo ou após uma grande atualização de versão, reserve 10 minutos para revisar estas configurações:\newline (1) Senha de bloqueio + biometria (2) Agrupamento de notificações + desativar push de marketing (3) Tamanho da fonte / Taxa de atualização / Modo escuro\newline (4) Teclado / Navegador / E-mail padrão (5) Contato de emergência / Ficha médica (6) Configurações de backup

\end{tcbinfo}

\booksection{V. Acessibilidade e Recursos Inclusivos: O Poder da Tecnologia para o Bem}

Não se assuste com o nome “Acessibilidade”. Essas funções são para todos. Foram originalmente projetadas para pessoas com deficiência, mas também são muito úteis para pessoas comuns — lesão temporária na mão, segurando um bebê, fone de ouvido quebrado no metrô, dificuldade para ler letras pequenas com óculos de leitura... tudo pode ser resolvido. Em última análise, o design acessível incorpora o conceito de\textbf{tecnologia para o bem}: permitir que todos possam desfrutar igualmente da vida digital.

\booksubsection{Os “Olhos” dos Usuários com Deficiência Visual — Guia Completo do Leitor de Tela}

O leitor de tela é a ferramenta central para usuários com deficiência visual usarem o smartphone. Ele\textbf{lê em voz alta todo o conteúdo visível na tela}, e o usuário opera o telefone por gestos e comandos de voz. Abaixo, apresentamos o uso do VoiceOver do iOS e do TalkBack do Android.

\begin{tcbinfo}{iOS VoiceOver}
\textbf{Caminho de ativação}: [Ajustes] > [Acessibilidade] > [VoiceOver], ative o interruptor. Também pode ser ativado/desativado rapidamente com três toques no botão lateral (ou botão de Início).

\textbf{Gestos básicos}:

\begin{enumerate}
  \item Deslizar com um dedo: navegar elemento por elemento pelo conteúdo da tela;
  \item Deslizar com dois dedos: rolar página inteira para cima/baixo e esquerda/direita;
  \item Toque duplo com um dedo: ativar o elemento em foco (equivalente a "tocar");
  \item Deslizar com três dedos: alternar páginas;
  \item Toque triplo com um dedo: ampliar ou selecionar (diferente do toque triplo comum).
\end{enumerate}

\textbf{Configurações personalizadas}: [VoiceOver] > [Fala], é possível ajustar a\textbf{velocidade de leitura}(padrão cerca de 50\%, pode ajustar para 70-80\% para uma audição mais fluida),\textbf{tom}e alternar entre diferentes pacotes de voz. O iOS possui “Voz aprimorada” e “Controlador de voz”: o primeiro reduz ruídos de respiração e pausas, tornando a audição mais natural.

\end{tcbinfo}

\begin{tcbinfo}{Android TalkBack}
\textbf{Caminho de ativação}: [Ajustes] > [Acessibilidade] > [TalkBack], ative o interruptor. Ou pressione e segure as teclas de volume para cima e para baixo por 3 segundos para ativação rápida (suportado por algumas marcas).

\textbf{Gestos básicos}:

\begin{enumerate}
  \item Deslizar um dedo para esquerda/direita: alternar entre elementos em foco;
  \item Toque duplo com um dedo: ativar elemento;
  \item Deslizar dois dedos para cima/baixo e esquerda/direita: rolar;
  \item Toque único com três dedos: entrar no modo de leitura (leitura contínua);
  \item Toque único com quatro dedos: pausar/continuar leitura.
\end{enumerate}

\textbf{Configurações personalizadas}: [TalkBack] > [Configurações de texto para fala], é possível ajustar a\textbf{velocidade},\textbf{tom}e alternar o\textbf{motor de voz}(recomenda-se instalar o motor de texto para fala do Google para uma leitura mais natural).

\end{tcbinfo}

\booksubsection{Aumentar Textos Pequenos, Clarear o Embaçado — Ajustes Visuais}

\booksubsection{Gesto de Ampliação}

\begin{itemize}
  \item \textbf{iOS}: [Ajustes] > [Acessibilidade] > [Zoom], após ativar,\textbf{toque a tela com três dedos}para ampliar. O padrão é zoom em tela cheia, com opção de zoom em janela (janela de ampliação local).\textbf{Afaste ou junte dois dedos para ajustar a ampliação (1,5x a 15x)}. Toque duplo com três dedos para travar a proporção de zoom atual.
  \item \textbf{Android}: [Ajustes] > [Acessibilidade] > [Ampliação]. Toque com três dedos para ampliar, toque novamente para sair; ou diretamente\textbf{toque três vezes na tela com três dedos}para ativar ampliação temporária (até 8x).
\end{itemize}

\booksubsection{Texto de Alto Contraste / Texto em Negrito}

\begin{itemize}
  \item \textbf{iOS}: [Ajustes] > [Acessibilidade] > [Tela e Tamanho do Texto], ative "Texto em Negrito", "Aumentar Contraste", "Reduzir Transparência", etc. "Inversão Inteligente" mantém o texto preto em fundo claro, evitando texto branco sobre fundo branco.
  \item \textbf{Android}: [Ajustes] > [Tela] > [Tamanho da fonte], ajustável em quatro níveis: "Pequeno, Normal, Grande, Muito Grande". Huawei/Xiaomi oferecem adicionalmente opções de "Negrito" e "Texto de alto contraste".
\end{itemize}

\booksubsection{Filtros de Cor (Amigável para Daltonismo)}

\begin{itemize}
  \item \textbf{iOS}: [Ajustes] > [Acessibilidade] > [Tela e Tamanho do Texto] > [Filtros de Cor], oferece\textbf{quatro modos para daltonismo: vermelho/verde/azul/amarelo}, além das opções "Escala de Cinza" e "Inverter Cores". Recomenda-se fazer primeiro um teste de visão de cores (integrado ao sistema) antes de selecionar o filtro correspondente.
  \item \textbf{Android}: [Ajustes] > [Acessibilidade] > [Correção de Cor], inclui\textbf{três modos para daltonismo}+ opção "Inverter cores".
\end{itemize}

\booksubsection{Descrição de Imagens por IA: Deixe a IA Ser Seus “Olhos”}

Quando o leitor de tela encontra uma imagem, a abordagem tradicional é pular ou apenas ler “imagem”. Mas agora a IA pode\textbf{descrever o conteúdo da imagem em uma frase}, sendo muito útil para reconhecer captchas, ler imagens e identificar menus.

\begin{tcbinfo}{Suporte à Descrição de Imagens por IA em Cada Sistema}
\begin{itemize}
  \item \textbf{iOS 15+}: Em [Ajustes] > [Acessibilidade] > [VoiceOver] > [Descrição de Fala], ative "Descrição de Imagens". O sistema usará um\textbf{modelo de IA local}para gerar automaticamente descrições de texto das imagens (requer iPhone 12 ou superior).
  \item \textbf{App TalkBack mais recente}: O Google TalkBack integra, nas versões mais recentes, a função de\textbf{perguntas sobre imagens}. É possível perguntar "O que há nesta imagem?" e a IA responde em tempo real (não depende estritamente da versão do Android, basta atualizar o app TalkBack para experimentar).
  \item \textbf{Apps de terceiros}: "Seeing AI" (iOS), "Lookout" (Android), "Xiao Yi Ajuda" (Huawei) são excelentes assistentes visuais com IA, capazes de identificar objetos, ler documentos e reconhecer expressões faciais.
\end{itemize}

\end{tcbinfo}

\begin{tcbpractice}{Cenários de Uso}
Tente, diante de uma foto de rua, perguntar “Quantas pessoas há nesta imagem? O que elas estão fazendo?”, ou abra um app de delivery e peça à IA para ler os nomes dos pratos no menu em imagem —\textbf{isso é verdadeiramente “tecnologia tornando a vida mais igualitária”}.

\end{tcbpractice}

\booksubsection{Os “Ouvidos” dos Usuários com Deficiência Auditiva — Guia Completo de Legendas em Tempo Real}

As Legendas em Tempo Real (Live Captions) são uma função de\textbf{conversão de fala em texto no nível do sistema}, extremamente úteis para usuários com deficiência auditiva em ambientes ruidosos, bibliotecas e situações com idiomas estrangeiros.

\begin{tcbinfo}{Ativação e Otimização das Legendas em Tempo Real}
\textbf{iOS}: [Ajustes] > [Acessibilidade] > [Legendas em Tempo Real], ative o interruptor. O sistema detectará todo o áudio do dispositivo (chamadas, vídeos, podcasts, FaceTime, memorandos de voz) e exibirá automaticamente as legendas na parte superior da tela.

\textbf{Android 12+}: Pressione a\textbf{tecla de volume}toque no ícone “Legendas em tempo real” na janela pop-up; ou ative em [Ajustes] > [Acessibilidade] > [Legendas em Tempo Real].

\textbf{Dicas de otimização da precisão}:

\begin{enumerate}
  \item Sempre que possível, use em\textbf{ambientes silenciosos}, longe de ar-condicionado e ruído de máquinas;
  \item A taxa de reconhecimento é maior quando o falante\textbf{articula claramente e fala em ritmo moderado};
  \item Diálogos com múltiplas pessoas reduzem a precisão. Use o modo "exibir apenas uma pessoa";
  \item Suporta mandarim, cantonês, inglês e outros idiomas principais.\textbf{A taxa de reconhecimento é menor para idiomas menos comuns};
  \item Se houver música de fundo ou interferência de vozes durante a fala, o reconhecimento será intermitente.
\end{enumerate}

\end{tcbinfo}

\booksubsection{Canal Duplo em Cenários de Chamada: Fala para Texto + Texto para Fala}

Quando o usuário\textbf{não consegue ouvir nem falar}(como perda súbita de voz, recuperação pós-cirurgia de garganta, gagueira severa), o telefone também pode atuar como “intérprete”.

\begin{tcbinfo}{Tradução de Chamadas em Tempo Real (TTY / RTT)}
\textbf{iOS}: [Ajustes] > [Acessibilidade] > [RTT/TTY]. Após ativar, durante a chamada, um teclado será exibido na tela. A fala do interlocutor é convertida em texto em tempo real, e\textbf{o texto que você digita será lido em voz alta}para o interlocutor.

\textbf{Android}: [Ajustes] > [Acessibilidade] > [Transcrição em Tempo Real] + [Notificação de Som]. “Transcrição em Tempo Real” exibe a fala do interlocutor como texto em tempo real; “Notificação de Som” converte sua entrada de texto em saída de voz quando não for conveniente falar.

\textbf{Compatibilidade}:

\begin{itemize}
  \item \textbf{RTT (Real-Time Text)}: Ambas as partes precisam ter suporte. O texto é exibido caractere por caractere conforme digitado, sem atraso. As principais operadoras já oferecem suporte;
  \item \textbf{TTY (Teletypewriter)}: Requer dispositivo TTY dedicado, sendo gradualmente substituído pelo RTT;
  \item A função de\textbf{conversão de voz em texto}do WeChat/QQ também é uma alternativa: pressione e segure a mensagem de voz "Converter para texto".
\end{itemize}

\end{tcbinfo}

\booksubsection{Funções “Pequenas mas Práticas” para Todos}

\begin{itemize}
  \item \textbf{Toque na Parte Traseira}(iOS 14+): [Ajustes] > [Acessibilidade] > [Toque] > [Toque na Parte Traseira], vincule\textbf{toque duplo/triplo}para acionar ações como captura de tela, bloqueio de tela, centro de controle, etc. Muito prático quando se está com as mãos ocupadas segurando encomendas.
  \item \textbf{Toque Assistivo / Botão Flutuante}: iOS [Acessibilidade] > [Toque] > [Toque Assistivo], Android [Acessibilidade] > [Navegação Flutuante]. Botão virtual que substitui gestos ou botões físicos. Salva a situação quando os botões não funcionam bem ou para operação com uma mão.
  \item \textbf{Controle por Interruptor}: Controla o telefone através de interruptores externos, toque único ou movimentos da cabeça.\textbf{Projetado para usuários com deficiência motora}. iOS [Acessibilidade] > [Controle por Interruptor], Android [Acessibilidade] > [Controle por Interruptor].
  \item \textbf{Reconhecimento de Sons}(iOS): [Acessibilidade] > [Reconhecimento de Sons]. Após ativar, o telefone "ouve" sons como campainha, choro de bebê, alarme de incêndio, sirenes, etc.\textbf{Ao reconhecer, vibra e exibe notificação na tela}. Muito valioso para usuários com deficiência auditiva e durante o sono.
  \item \textbf{Alerta de LED Piscante}: iOS [Acessibilidade] > [Áudio/Visual] > [LED Piscante para Alertas], Android [Ajustes] > [Notificações] > [Alerta de Flash]. O flash pisca ao receber chamadas e notificações.\textbf{Permite perceber alertas mesmo em ambientes ruidosos sem som}.
  \item \textbf{Áudio Mono}: [Ajustes] > [Acessibilidade] > [Áudio] > [Áudio Mono]. Combina os canais esquerdo e direito. Usuários com perda auditiva em um ouvido não perdem a metade da informação sonora.
\end{itemize}

\begin{tcbinfo}{A Filosofia de Design Acessível da Apple}
A Apple investe enormemente em acessibilidade, atualizando e adicionando novos recursos de acessibilidade todos os anos. O site oficial da Apple possui uma página dedicada à acessibilidade, detalhando o uso de cada função. Os recursos de acessibilidade do iOS costumam ser mais completos e a experiência mais fluida do que no Android. Esta é uma razão importante pela qual muitos usuários com deficiência escolhem a Apple.

\end{tcbinfo}

\begin{tcbpractice}{Experimente por 30 segundos}
Reserve um momento livre para experimentar o seguinte:

\begin{enumerate}
  \item Ative o\textbf{VoiceOver / TalkBack}e tente navegar pelo WeChat deslizando com três dedos sem usar os olhos;
  \item Ative as\textbf{Legendas em Tempo Real}e fale para o telefone para ver a precisão do reconhecimento;
  \item Use o\textbf{toque com três dedos}para ampliar um texto em 200\% e experimente o cotidiano de um usuário com deficiência visual;
  \item Desative todas as funções anteriores e reflita: "Se eu precisasse dessas funções para sempre, como eu gostaria que fossem projetadas?".
\end{enumerate}

\end{tcbpractice}
